\section{Terraform}
\textbf{IaC:} manage \& provision infrastructure (\kf{compute, storage, network}) through code instead of manual processes (Red Hat). \\
\textbf{Motivation:} replace slow, drift-prone \kf{manual} console clicking with code; Automation pays off but carries its own upkeep cost. \\
\textbf{Benefits:} \kf{idempotent} (declarative, re-apply = same result), \kf{drift detection} (state vs.\ reality), version-controlled (history/review/rollback), \kf{consistency} (change a source $\to$ all references update), \kf{reusable \& standardized} (modules: write once, deploy many, enforce best practices), \kf{team collaboration} (shared, locked remote state), \kf{flexible} (variables, no hard-coded secrets). \\
\textbf{Declarative vs. Imperative:} Terraform is \kf{declarative} -- you describe the desired end state and TF figures out the steps to get there. \kf{Imperative} (shell scripts, manual CLI calls) instead describes \kf{how} to reach a result step by step. \\
\textbf{Toolchain (4 components):} \kf{language} (HCL -- declare infra \& vars), \kf{state} (\texttt{.tfstate}, maps config to real resources), \kf{provisioner} = Terraform-core plugins: Providers (access a cloud's API: AWS/Azure/GCP/K8s) + Provisioners (\texttt{remote-exec}/\texttt{local-exec} scripts on a resource), \kf{backend} (where state lives: local or remote S3+lock; \texttt{init} sets it up + downloads plugins). \\
TF doesn't speak directly with an SDK, but rather Terraform -> Provider -> Client SDK. Diffrent providers enable diffrent platforms (AWS, Azure, Kubernetes, ...). \textbf{Resources} (\texttt{resource}) are the "what" TF creates and manages (e.g. an EC2 instance). \textbf{Data Sources} (\texttt{data}) are read-only lookups of existing infrastructure used as input (e.g. an existing AMI or VPC). A sample in HCL (Hashicorp Configuration Language):
\begin{lstlisting}[language=hcl]
terraform {
  required_version = ">= 1.0.0"
  required_providers { aws = { source = "hashicorp/aws", version = "~> 5.0" } }
}
provider "aws" { region = "us-east-1" }
variable "instance_type" { default = "t2.micro" }
resource "aws_instance" "web" {
  ami           = data.aws_ami.ubuntu.id
  instance_type = var.instance_type
}
output "public_ip" { value = aws_instance.web.public_ip }
\end{lstlisting}
To deploy infrastructure, write HCL in files like \texttt{main.tf}, then run \texttt{terraform init}, \texttt{terraform plan} to show changes that would be made, \texttt{terraform apply} to actually apply the changes. Use \texttt{terraform destroy} to delete all made changes.
\subsection{OpenTofu}
Fork of Terraform v1.5.7 under \textbf{Linux Foundation}, MPL-2.0 (TF switched to BUSL). Drop-in replacement, adds \textbf{state encryption}. Avoids single-vendor lock-in.
\subsection{HCL Features}
\textbf{count} (numeric loop, \texttt{count.index}) and \textbf{for\_each} (map/set loop, \texttt{each.key/value}) for multiple instances. \textbf{Data source} example:
\begin{lstlisting}[language=hcl]
data "aws_ami" "ubuntu" {
  most_recent = true
  filter { name = "name"   values = ["ubuntu/images/hvm-ssd/ubuntu-*"] }
  filter { name = "virtualization-type" values = ["hvm"] }
  owners = ["099720109477"] # Canonical
}
\end{lstlisting}
\subsection{Variables}
Precedence: CLI \texttt{-var}/\texttt{-var-file} > env \texttt{TF\_VAR\_<name>} > \texttt{*.tfvars} > default. No default \& not supplied -> CLI prompts operator. \texttt{TF\_VAR\_} prefix ideal for secrets (\texttt{TF\_VAR\_access\_key=...}). \texttt{depends\_on} forces explicit ordering; implicit deps come from attribute references.
\subsection{State}
\texttt{terraform.tfstate} (JSON) maps HCL to real resource IDs (e.g. ARN). \textbf{3-way diff} on apply: \textbf{NEW} (HCL) vs \textbf{PREVIOUS} (.tfstate) vs \textbf{EXISTING} (cloud refresh) -> decides actions, detects drift. \textbf{Remote backend} (S3) shares state in teams; \textbf{state locking} via \texttt{use\_lockfile} (S3-native, replaces DynamoDB):
\begin{lstlisting}[language=hcl]
terraform {
  backend "s3" {
    bucket       = "my-tfstate"
    key          = "prod/terraform.tfstate"
    region       = "us-east-1"
    use_lockfile = true
  }
}
\end{lstlisting}
\subsection{File \& Folder Layout}
\treeline{\texttt{├─~}\kf{\texttt{modules/}}}
\treeline{\texttt{│~~~~~~}shared modules, reusable across envs}
\treeline{\texttt{├─~}\kf{\texttt{dev/}}}
\treeline{\texttt{│~~├──~}\kf{\texttt{terraform.tf}}}
\treeline{\texttt{│~~│~~~}required\_version + providers}
\treeline{\texttt{│~~├──~}\kf{\texttt{backend.tf}}}
\treeline{\texttt{│~~│~~~}remote state (kept out of VCS)}
\treeline{\texttt{│~~├──~}\kf{\texttt{providers.tf}}}
\treeline{\texttt{│~~│~~~}provider auth/config}
\treeline{\texttt{│~~├──~}\kf{\texttt{main.tf}}}
\treeline{\texttt{│~~│~~~}resources \& data sources}
\treeline{\texttt{│~~├──~}\kf{\texttt{variables.tf}}}
\treeline{\texttt{│~~│~~~}inputs (alphabetical)}
\treeline{\texttt{│~~├──~}\kf{\texttt{outputs.tf}}}
\treeline{\texttt{│~~│~~~}returns (alphabetical)}
\treeline{\texttt{│~~├──~}\kf{\texttt{locals.tf}}}
\treeline{\texttt{│~~│~~~}derived values}
\treeline{\texttt{│~~└──~}\kf{\texttt{override.tf}}}
\treeline{\texttt{│~~~~~~}loaded last (use sparingly)}
\treeline{\texttt{└──~}\kf{\texttt{prod/}}}
\treeline{\texttt{~~~~~~~}same layout, its own state}
\smallskip
Per-env folders $\to$ separate state, so a corruption hits only one env.
\subsection{Modules}
Reusable folder of \texttt{.tf} files (root config is itself a module):
\begin{lstlisting}[language=hcl]
module "vpc" {
  source = "./modules/vpc" # or registry/git URL
  cidr   = "10.0.0.0/16"   # input variable
} # reference an output: module.vpc.vpc_id
\end{lstlisting}
\subsection{Patterns \& Recovery}
\texttt{import} adopts existing infra; \texttt{terraform apply -replace=<addr>} (or legacy \texttt{taint}) forces a destroy+recreate of a single resource. Region-aware AMI via map:
\begin{lstlisting}[language=hcl]
import { to = aws_instance.web   id = "i-0abc123" }   # then: apply

variable "AMIS" { type = map(string)  default = { us-east-1 = "ami-05b1..." } }
ami = var.AMIS[var.region]
\end{lstlisting}
